\section{Próximo gate discriminante}
O próximo investimento deve começar por source/static feasibility, não por uma nova campanha física longa. C189 propõe um screen limitado para perceber se \texttt{ubatch64} no caminho C75/waves pode reduzir repetição de prefill para o nominal153, tratando-o como \textbf{novo perfil numérico} devido às divergências shape-dependent já observadas em C4--C6. \evidence{E-S01-C189}

A ordem recomendada é:
\begin{enumerate}
\item localizar no source/trace uma repetição ou accounting evitável que ub64 ataque;
\item se existir, congelar identidade de modelo/backend/workload e executar gates numéricos/funcionais próprios;
\item apenas depois fazer um screen pareado e limitado de first-final, protegendo decode e recursos;
\item parar se o source screen não encontrar trabalho removível ou se correctness divergir sem uma qualificação explícita.
\end{enumerate}

Em paralelo, qualquer claim de I/O deve medir o tier certo. \texttt{read\_bytes}, bytes de tensor e requests de copy não substituem instrumentação de tráfego físico. Amdahl/Roofline servem para limitar o ganho possível antes de acrescentar mecanismos \cite{amdahl1967validity,williams2009roofline}.
